% cSpell:ignore VLSI,NMOS,PMOS,FDSOI,CMOS,TTL,OD,TG
\chapter{门电路}
\setlength{\parskip}{0.4em}
\setlength{\itemsep}{0.2em}
\setlength{\tabcolsep}{4pt} % 统一表格列间距

\section{概述}
\subsection{数字系统核心思想}
数字系统将连续模拟量抽象为离散0/1二值逻辑，核心目标是可靠区分逻辑0、1，依靠统一电平标准提升电路可靠性。

\subsubsection{模拟系统缺陷}
模拟信号连续无明确分界，存在三大问题。噪声直接造成信号失真。复制，取反运算无法无误差实现。对器件精度，电源稳定性要求高，抗干扰差。

\subsubsection{数字系统优缺点}
\begin{itemize}
    \item 优势：输出标准0/1电平，可容忍带噪声的劣质输入。统一接口便于模块化。高低电平存在容错区间，降低硬件要求。
    \item 代价：模拟信号数字化存在量化损失。需要额外模数转换电路。
\end{itemize}

\subsection{门电路定义}
实现与、或、非及与非、或非、异或等复合逻辑的基础单元电路。

\subsection{正负逻辑（本书采用正逻辑）}
\begin{table}[htbp]
    \centering
    \caption{正负逻辑电平映射}
    \begin{tabular}{|c|c|c|}
        \hline
        逻辑类型 & 高电平 & 低电平 \\
        \hline
        正逻辑 & 逻辑1 & 逻辑0 \\
        \hline
        负逻辑 & 逻辑0 & 逻辑1 \\
        \hline
    \end{tabular}
\end{table}

\subsection{高低电平实现原理}
用开关通断实现两种电平。开关断开输出\(V_{CC}\)（逻辑1），开关闭合输出0V（逻辑0）。实际电路依靠二极管，三极管截止/饱和等效开关。

\subsection{数字电路分类与集成规模}
\subsubsection{优点}
高低电平具备容错电压区间，抗干扰强，易大规模集成，对器件精度要求低。

\subsubsection{分类}
1. 分立元件门：电阻，电容，半导体器件独立搭建。
2. 集成逻辑门：器件集成于单片半导体芯片。

\begin{table}[htbp]
    \centering
    \caption{集成电路规模分级}
    \begin{tabular}{|c|c|}
        \hline
        缩写 & 单片器件数量 \\
        \hline
        SSI 小规模 & $\leq 100$ \\
        \hline
        MSI 中规模 & $100 \sim 1000$ \\
        \hline
        LSI 大规模 & $10^3 \sim 10^5$ \\
        \hline
        VLSI 超大规模 & $> 10^5$ \\
        \hline
    \end{tabular}
\end{table}

\subsection{本章内容}
依次讲解二极管门电路，CMOS门电路，TTL门电路。

\section{二极管门电路}
\subsection{二极管开关等效模型}
数字电路使用3种等效模型。
1. 完整模型：包含导通电阻\(r_D\)，导通压降\(V_{ON}\)。
2. 简化模型：仅保留导通压降，忽略导通内阻。
3. 理想模型：导通压降、内阻均为0，数字电路首选。

\subsection{二极管稳态开关特性}
依靠单向导电性，输入电平控制通断。
- \(v_I=V_{CC}\)（高电平）：二极管截止，输出\(v_O=V_{CC}\)（逻辑1）。
- \(v_I=0\mathrm{V}\)（低电平）：二极管导通，输出\(v_O=0\mathrm{V}\)（逻辑0）。

\subsection{二极管动态开关特性}
电压突变时二极管切换存在延时，根源为PN结存储电荷。
1. 开通时间\(t_{on}\)：截止→导通，建立内电场，时间极短。
2. 反向恢复时间\(t_{re}\)：导通→截止，消散存储电荷，\(\boldsymbol{t_{re} \gg t_{on}}\)。

低频电路可忽略延时，高频电路反向恢复时间不可忽略。

\subsection{二极管与门}
逻辑式：\(\boldsymbol{Y=A\cdot B}\)，仅输入全高电平时输出高电平，其余输入组合输出低电平。
缺陷：输出电平偏移，带负载能力差，无法实现非逻辑。

\subsection{二极管或门}
逻辑式：\(\boldsymbol{Y=A+B}\)，任意一路输入高电平则输出高电平，仅全低输入时输出低电平。

\subsection{二极管门电路通用缺陷}
1. 输出电平存在偏移，级联后电平失真累积。
2. 带负载能力弱。
3. 无取反功能，仅能用于芯片内部简易电路。

\section{CMOS门电路}
\subsection{MOS管开关特性（重点：两类沟道导通条件）}
CMOS电路由**NMOS，PMOS互补配对**构成，二者导通条件完全相反。

\textbf{1. 增强型NMOS（N沟道）导通条件}
栅源电压大于开启阈值电压才能导通：
\[
\boldsymbol{V_{GS} \ge V_{TH} \implies \text{管子导通}}
\]
若\(V_{GS}<V_{TH}\)，管子截止。

\textbf{2. 增强型PMOS（P沟道）导通条件}
栅源电压小于负阈值电压才能导通：
\[
\boldsymbol{V_{GS} \le -|V_{TH}| \implies \text{管子导通}}
\]
若\(V_{GS}>-|V_{TH}|\)，管子截止。

MOS管四端定义：源极S，栅极G，漏极D，衬底B，器件物理对称。
电路布线规则：NMOS漏极接高电位，PMOS源极接高电位。

\subsection{MOS管工艺演进说明}
FinFET工艺普及后，工艺纳米数值不再代表栅极/Fin实际物理宽度，仅为摩尔定律代际命名。例如7nm制程无任何结构尺寸等于7nm。
演进路线：平面体硅CMOS → FinFET鳍式MOS → FDSOI绝缘硅。

\subsection{CMOS反相器结构与工作原理}
CMOS全称互补MOS，反相器由PMOS管\(T_1\)，NMOS管\(T_2\)构成互补对称电路，输入同时驱动两管栅极，两管始终一导通一截止。

设两管开启电压\(V_{GS(th)P}\)，\(V_{GS(th)N}\)，且\(V_{DD}>|V_{GS(th)P}|+V_{GS(th)N}\)。
1. 输入\(v_I=V_{IL}=0\mathrm{V}\)低电平：\(T_1(\text{PMOS})\)导通，\(T_2(\text{NMOS})\)截止，输出\(v_O=V_{DD}\)（逻辑1）。
2. 输入\(v_I=V_{IH}=V_{DD}\)高电平：\(T_1(\text{PMOS})\)截止，\(T_2(\text{NMOS})\)导通，输出\(v_O=0\mathrm{V}\)（逻辑0）。

逻辑功能：\(\boldsymbol{Y=\overline{A}}\)。

\subsection{CMOS反相器电压传输特性}
描述输入\(v_I\)与输出\(v_O\)的对应关系，分为三段。
1. AB段（输入低压区，\(v_I < V_{GS(th)N}\)）：\(T_1\)导通，\(T_2\)截止，输出维持高电平\(V_{OH}=V_{DD}\)。
2. BC段（过渡区，\(V_{GS(th)N} < v_I < V_{DD}-|V_{GS(th)P}|\)）：\(T_1\)，\(T_2\)同时导通，输入微小变化即可使输出快速翻转，翻转中点为\(V_{DD}/2\)。此区间电源到地存在持续电流，功耗大，电路不可长期工作在此区间。
3. CD段（输入高压区，\(v_I > V_{DD}-|V_{GS(th)P}|\)）：\(T_1\)截止，\(T_2\)导通，输出维持低电平\(V_{OL}=0\)。

\subsection{CMOS反相器电流传输特性}
描述漏极电流\(i_D\)随输入\(v_I\)变化规律，分段与电压传输特性对应。
1. AB，CD段：单管导通，另一管截止，漏极电流近似为0。
2. BC过渡区：两管同时导通，漏极出现峰值电流，在\(v_I=V_{DD}/2\)处电流达到最大值。

\subsection{噪声容限}
噪声容限是保证输出高低电平不变时，输入允许叠加干扰的最大幅值，分为高电平噪声容限\(V_{NH}\)，低电平噪声容限\(V_{NL}\)。
\[
\begin{cases}
V_{NH}=V_{OH(\text{min})}-V_{IH(\text{min})}\\
V_{NL}=V_{IL(\text{max})}-V_{OL(\text{max})}
\end{cases}
\]

提高电源电压\(V_{DD}\)可同步提升\(V_{NH}\)，\(V_{NL}\)，增强抗干扰能力。直流噪声容限随输入脉冲宽度变化，窄脉冲下噪声容限更高。

\subsection{CMOS反相器静态输入输出特性}
\subsubsection{输入特性与输入保护电路}
MOS栅极与衬底间为薄层\(\mathrm{SiO_2}\)绝缘层，易被静电击穿，必须增设输入保护电路。

保护电路由钳位二极管\(\mathrm{D_1}\)，\(\mathrm{D_2}\)与限流电阻\(R_S\)构成。
- 正常输入范围\(0 \le v_I \le V_{DD}\)，二极管截止，保护电路无作用。
- \(v_I > V_{DD}+V_F\)时\(\mathrm{D_1}\)导通，栅极电位被钳位。\(v_I < -V_F\)时\(\mathrm{D_2}\)导通，限制栅极负压，避免栅电容过压击穿。

输入特性曲线：正常区间输入电流近似为0。输入超出\([-0.7\mathrm{V},V_{DD}+0.7\mathrm{V}]\)时二极管导通，输入电流急剧增大。

\subsubsection{输出特性}
输出特性描述输出电压与输出负载电流的关系，分高低电平两种工况。
1. 低电平输出（输入高电平\(v_I=V_{DD}\)）：\(T_1\)截止，\(T_2\)导通，负载电流灌入\(T_2\)。负载电流增大时，输出低电平\(V_{OL}\)抬升。
2. 高电平输出（输入低电平\(v_I=0\)）：\(T_1\)导通，\(T_2\)截止，电流从\(T_1\)流向负载。负载电流增大时，输出高电平\(V_{OH}\)下降。

\subsection{CMOS反相器动态特性}
\subsubsection{传输延迟时间}
负载电容\(C_L\)充放电是产生传输延迟的核心原因，定义两类延迟。
- \(t_{PHL}\)：输入上升沿50\%幅值到输出下降沿50\%幅值的时间。
- \(t_{PLH}\)：输入下降沿50\%幅值到输出上升沿50\%幅值的时间。

延迟大小与负载电容\(C_L\)，电源电压\(V_{DD}\)相关。

\subsubsection{动态功耗}
CMOS静态功耗极低，可忽略。总动态功耗分为两部分。
1. 导通功耗\(P_T\)：电平切换过渡区间两管同时导通产生瞬时电流，平均功耗：
\[
P_T = V_{DD} I_{TAV},\quad I_{TAV}=\frac{1}{T}\left(\int_{t_1}^{t_2}i_T dt + \int_{t_3}^{t_4}i_T dt\right)
\]
2. 负载电容充放电功耗\(P_C\)：负载电容充放电产生，公式：
\[
P_C = C_L f V_{DD}^2
\]
总动态功耗：\(\boldsymbol{P_D = P_T + P_C}\)。

\subsubsection{扇出特性}
扇出代表门电路带负载能力。
1. 静态工况：CMOS输入电阻极高（\(10^{10}\sim10^{12}\Omega\)），输出导通内阻低，低频（\(<1\mathrm{MHz}\)）下扇出可达50以上。
2. 动态工况：负载门等效输入电容随负载数量增加而变大，电容充放电拉长电平切换时间，频率越高，可驱动负载数量越少。

\subsection{CMOS反相器核心特点}
1. 静态功耗极低：高低电平稳态下总有一管截止，静态电流仅微安级，属于微功耗电路。
2. 传输速度快：导通管内阻小，对容性负载充放电速度快，平均传输延迟\(t_{pd}\approx10\mathrm{ns}\)。
3. 噪声容限高：阈值转换点位于\(V_{DD}/2\)，噪声容限最高可接近电源电压的50\%，抗干扰能力强。
4. 静态扇出大，动态扇出受工作频率与负载电容约束。

\subsection{其他类型CMOS逻辑门}
\subsubsection{CMOS与非门、或非门}
1. 二输入与非门：PMOS\(T_1\)，\(T_3\)并联，NMOS\(T_2\)，\(T_4\)串联。逻辑式\(Y=(AB)'\)。任意输入为0，PMOS支路导通，NMOS支路断开，输出高电平。全输入为1，NMOS全部导通，输出低电平。
2. 二输入或非门：PMOS\(T_1\)，\(T_3\)串联，NMOS\(T_2\)，\(T_4\)并联。逻辑式\(Y=(A+B)'\)。任意输入为1，NMOS支路导通，PMOS支路断开，输出低电平。全输入为0，PMOS全部导通，输出高电平。

普通与非/或非门缺陷：输出导通电阻随输入组合变化，高低电平受输入端数量影响。
解决方案：输入端，输出端各加一级反相缓冲器，构成带缓冲级CMOS门，输出电阻，高低电平不受输入状态影响，电压传输特性更陡。

\subsubsection{漏极开路输出门（OD门）}
输出级仅保留NMOS漏极开路，使用时必须外接上拉电阻\(R_L\)至电源。
1. 功能：实现电平转换，大负载驱动，多路输出“线与”。
2. 线与逻辑：多个OD门输出直接并联，输出结果为各门输出的与运算。逻辑式\(Y=(AB)'(CD)'=(AB+CD)'\)。
3. 上拉电阻取值约束：
\[
\begin{cases}
R_{L(\text{max})}=\dfrac{V_{DD}'-V_{OH(\text{min})}}{nI_{OH}+mI_{IH}}\\[6pt]
R_{L(\text{min})}=\dfrac{V_{DD}'-V_{OL(\text{max})}}{I_{OL(\text{max})}-m'I_{IL}}
\end{cases}
\]
\(n\)为并联OD门数量，\(m\)，\(m'\)为高低电平负载门数量，取值需兼顾高低电平标准。

\subsubsection{CMOS传输门（TG）【重点：传输导通/截止条件】}
由一对互补NMOS，PMOS并联构成，互补控制信号\(C\)，\(C'\)驱动，双向导通器件。

\textbf{1. 导通条件（传输模式）}
控制电平：\(\boldsymbol{C=1,\ C'=0}\)；
输入电压范围：\(\boldsymbol{0 < v_I < V_{DD}}\)；
此时NMOS，PMOS同时导通，输入信号完整传递至输出：\(v_O=v_I\)。

\textbf{2. 截止条件（高阻隔离）}
控制电平：\(\boldsymbol{C=0,\ C'=1}\)；
NMOS，PMOS同时截止，输出端与输入完全电气隔离，呈现高阻态。

\textbf{3. 关键特性}
- 双向传输：输入，输出可互换，支持连续模拟电压；
- 配套电路：与反相器组合构成通用双向模拟开关；
- 传输系数：\(K_{TG}=\dfrac{R_L}{R_{TG}+R_L}\)，负载电阻\(R_L \gg R_{TG}\)时传输系数趋近1，信号衰减极小。

\subsubsection{三态输出CMOS门（三态门）}
输出具备高电平，低电平，高阻三种状态，带低电平有效使能端\(EN'\)。
1. 工作逻辑：
    - \(EN'=0\)：电路正常工作，\(Y=A'\)。
    - \(EN'=1\)：输出高阻态，与后级电路电气隔离。
2. 典型应用：
    - 总线结构：多个三态门并联至同一总线，分时控制使能端实现多路数据传输，任意时刻仅允许一个门使能。
    - 双向数据传输：两组三态门反向对接，分时导通实现总线双向收发数据。

\section{TTL门电路}
\subsection{双极型三极管（BJT）开关基础}
TTL电路基于NPN双极三极管实现开关逻辑，硅管导通开启电压\(V_{ON}\approx0.7\mathrm{V}\)。

\subsubsection{三极管三个工作区}
1. 截止区：\(V_{BE}<V_{ON}\)，\(i_B\approx0\)，\(i_C\approx0\)，集射极近似开路。
2. 放大区：\(V_{BE}>V_{ON}\)，\(V_{CE}>0.7\mathrm{V}\)，\(i_C=\beta i_B\)，电流线性放大。
3. 饱和区：\(V_{CE}<0.7\mathrm{V}\)，\(i_C\)不再随\(i_B\)线性增大，饱和压降\(V_{CE(sat)}\approx0.1\sim0.3\mathrm{V}\)，集射极近似短路。

\subsubsection{三极管静态开关电路}
输入低电平\(v_I=V_{IL}<V_{ON}\)，三极管截止，输出\(v_O=V_{CC}\)高电平。
输入高电平\(v_I=V_{IH}\)，基极电流足够使管子饱和，输出\(v_O=V_{CE(sat)}\)低电平，实现反相功能\(Y=\overline{A}\)。

临界饱和基极电流：
\[
I_{BS}=\frac{V_{CC}-V_{CE(sat)}}{\beta R_C}
\]
满足\(i_B>I_{BS}\)三极管可靠饱和。

\subsubsection{动态开关特性}
三极管导通/关断存在电荷建立，消散过程，输出电压变化滞后输入，产生开关延迟。

\subsection{TTL反相器电路结构与工作原理}
标准TTL反相器由三级构成：输入级，倒相级，推拉式输出级，电源\(V_{CC}=5\mathrm{V}\)。
1. 输入级：\(T_1\)与钳位二极管\(D_1\)，抑制负向输入脉冲，保护发射结。
2. 倒相级：\(T_2\)，实现信号相位翻转。
3. 输出级：\(T_4\)，\(T_5\)构成推拉输出，任意时刻一管导通一管截止，降低静态功耗，提升带载能力。二极管\(D_2\)保证\(T_4\)可靠截止。

\subsubsection{电压传输特性分段}
1. AB段（截止区，\(v_I<0.6\mathrm{V}\)）：\(T_2\)，\(T_5\)截止，\(T_4\)导通，输出高电平\(V_{OH}\approx3.4\mathrm{V}\)。
2. BC段（线性区，\(0.7\mathrm{V}<v_I<1.3\mathrm{V}\)）：\(T_2\)放大导通，\(T_4\)导通，输入上升，输出线性下降。
3. CD段（转折区，\(v_I\approx V_{TH}=1.4\mathrm{V}\)）：\(T_2\)，\(T_5\)同时导通，\(T_4\)截止，输出快速跌落至低电平。
4. DE段（饱和区，\(v_I\)继续升高）：\(T_5\)深度饱和，输出维持\(V_{OL}\approx0.3\mathrm{V}\)不变。

\subsubsection{电路核心特点}
1. \(T_1\)倒置工作，电流放大系数小于1。
2. 推拉输出结构：高电平由\(T_4\)输出电流，低电平由\(T_5\)灌入电流，输出电阻均较小，带载能力强。
3. 输入钳位二极管\(D_1\)限制负输入电流，保护输入管。

\subsection{TTL反相器静态输入、输出特性}
\subsubsection{输入特性}
1. 输入低电平\(v_I=0.2\mathrm{V}\)：产生灌入输入的低电平电流\(I_{IL}\approx-1\mathrm{mA}\)。输入短路电流\(I_{IS}\approx-1.075\mathrm{mA}\)。
2. 输入高电平\(v_I=3.4\mathrm{V}\)：\(T_1\)发射结反偏倒置，仅有微小漏电流\(I_{IH}<40\mu\mathrm{A}\)。

\subsubsection{输入端负载特性}
输入端外接电阻\(R_P\)接地时，\(R_P\)压降即为输入电压\(v_I\)：
\[
v_I=\frac{R_P}{R_1+R_P}(V_{CC}-v_{BE1})
\]
- \(R_P \leq 0.7\mathrm{k}\Omega\)：输入为低电平。
- \(R_P \geq 1.5\mathrm{k}\Omega\)：输入为高电平。

\(v_I\)最高被钳位在约1.4V，不再随\(R_P\)增大继续上升。

\subsubsection{输出特性}
1. 高电平输出：\(T_4\)导通，\(T_5\)截止，输出电流拉出，负载电流增大时\(V_{OH}\)小幅下降，74系列最大拉出电流\(|I_{OH(max)}|=0.4\mathrm{mA}\)。
2. 低电平输出：\(T_5\)饱和导通，\(T_4\)截止，负载电流灌入，电流越大\(V_{OL}\)略微抬升，74系列最大灌电流\(I_{OL(max)}=16\mathrm{mA}\)。

\subsubsection{扇出系数计算}
扇出\(N\)为门电路可驱动同系列负载门最大数量，分别计算高低电平可驱动数量后取最小值：
\[
\begin{cases}
N_1 \leq \dfrac{I_{OL(\text{max})}}{|I_{IL}|} \quad (\text{低电平约束})\\[4pt]
N_2 \leq \dfrac{I_{OH(\text{max})}}{I_{IH}} \quad (\text{高电平约束})
\end{cases}
\]
最终\(N=\min(N_1,N_2)\)。以74系列参数为例：\(N_1=16\)，\(N_2=10\)，最终扇出\(N=10\)。

\subsection{TTL反相器动态特性}
\subsubsection{传输延迟时间}
平均传输延迟\(t_{pd}\)衡量工作速度，分为\(t_{PHL}\)（输入上升沿到输出下降沿延迟），\(t_{PLH}\)（输入下降沿到输出上升沿延迟），由结电容，寄生电容充放电导致，74系列典型延时3~40ns。

\subsubsection{动态尖峰电流与交流噪声容限}
1. 电平切换瞬间\(T_4\)，\(T_5\)短暂同时导通，产生电源尖峰电流，会通过电源线形成系统干扰。
2. 窄脉冲输入时输出变化滞后，交流噪声容限远大于直流噪声容限，脉冲宽度越窄，噪声容限越高。

\subsection{其他类型TTL逻辑门}
\subsubsection{TTL与非门、或非门、与或非门、异或门}
1. 与非门：多发射极输入管实现与逻辑，后级和反相器结构一致，逻辑式\(Y=(AB)'\)，输入高低电平电流随并联输入端数量成倍变化。
2. 或非门：多个输入级并联实现或逻辑，逻辑式\(Y=(A+B)'\)。
3. 与或非门：多级与输入再做或非运算，逻辑式\(Y=(AB+CD)'\)。
4. 异或门：由与非门组合实现，逻辑式\(Y=A\oplus B=A'B+AB'\)。

\subsubsection{OC集电极开路门}
普通推拉式TTL输出不可直接并联，会出现大电流烧坏器件。去掉上拉推拉管得到OC门。
1. 特点：输出管耐压，耐流能力强，外接独立上拉电阻可实现电平转换。
2. 线与功能：多个OC输出并联，实现输出逻辑与：\(Y=(AB)'(CD)'=(AB+CD)'\)。
3. 上拉电阻\(R_L\)计算和CMOS OD门公式一致，需同时满足高低电平规范。

\subsubsection{TTL三态门（TSL）}
在普通与非门增加使能控制端，具备高，低电平，高阻三态。
1. 使能有效时正常执行与非逻辑，使能无效时输出高阻。
2. 用于总线分时传输，双向数据收发，同一总线同一时刻仅允许一个三态门使能。

\subsection{TTL与CMOS电路接口规范}
两类电路互连时，驱动端必须满足负载端电平，电流规范：
\[
\begin{cases}
V_{OH(\text{min})} \ge V_{IH(\text{min})}\\
V_{OL(\text{max})} \le V_{IL(\text{max})}\\
I_{OH(\text{max})} \ge nI_{IH(\text{max})}\\
I_{OL(\text{max})} \ge m|I_{IL(\text{max})}|
\end{cases}
\]

\subsubsection{TTL驱动CMOS}
- 低压CMOS场景：TTL输出增加上拉电阻提升高电平，使\(V_{OH}\approx V_{DD}\)。
- 高压CMOS场景：改用OC门抬高输出电平。

\subsubsection{CMOS驱动TTL}
- 电流不足时：可将多个CMOS门并联增大灌电流，或后端增加专用驱动芯片提升带载能力。